\section{The complete D-Tracker model}
\label{sec:model}

\subsection{Three components added together}

For the current window $\mathcal{X}^{c}$, the paper assumes that observations are approximately the
sum of a trend tensor, a seasonal tensor, and an outlier tensor (paper Eq.~1):
\begin{equation}
  \mathcal{X}^{c}\;\approx\;\widehat{\mathcal{X}}_d
  +\widehat{\mathcal{X}}_s+\widehat{\mathcal{X}}_o,
  \qquad
  \widehat{\mathcal{X}}_d,\widehat{\mathcal{X}}_s,\widehat{\mathcal{X}}_o
  \in\mathbb{R}^{k\times l\times L_c}.
  \label{eq:model-sum}
\end{equation}
The decomposition is not decoration. Each component carries a distinct assumption:
\begin{itemize}
  \item $\widehat{\mathcal{X}}_d$ contains \textbf{trend and diffusion}, the structured changing part;
  \item $\widehat{\mathcal{X}}_s$ contains \textbf{seasonality}, the repeating part;
  \item $\widehat{\mathcal{X}}_o$ contains \textbf{outliers}, assumed sparse and unusual.
\end{itemize}
Without the split, a promotion could be learned as a long-term trend; without the outlier component,
one shock could corrupt every coefficient.

\subsection{Trend: compress a core and project it back}

The trend tensor still has size $k\times l\times L_c$, too large to read directly. The paper
factorizes it into a small core and two factor matrices (paper Eq.~2):
\begin{equation}
  \widehat{\mathcal{X}}_d
  =\mathcal{W}^{(\mathrm{core})}
   \times_{\mathrm{key}}\mathbf{W}^{(\mathrm{key})}
   \times_{\mathrm{loc}}\mathbf{W}^{(\mathrm{loc})},
  \label{eq:trend-factorization}
\end{equation}
with
\[
  \mathcal{W}^{(\mathrm{core})}\in\mathbb{R}^{d_k\times d_l\times L_c},
  \quad
  \mathbf{W}^{(\mathrm{key})}\in\mathbb{R}_{+}^{d_k\times k},
  \quad
  \mathbf{W}^{(\mathrm{loc})}\in\mathbb{R}_{+}^{d_l\times l},
\]
and $d_k\ll k$, $d_l\ll l$. Elementwise, this is
\begin{equation}
  \widehat{x}_{d,uv}(t)
  =\sum_{i=1}^{d_k}\sum_{j=1}^{d_l}
    w^{(\mathrm{core})}_{ij}(t)\,
    W^{(\mathrm{key})}_{iu}W^{(\mathrm{loc})}_{jv}.
  \label{eq:trend-elementwise}
\end{equation}
Each symbol has a concrete meaning: $u$ is a real keyword, $v$ a real location, and $i,j$ latent
groups. If $W^{(\mathrm{key})}_{iu}$ is large, keyword $u$ ``belongs'' to group $i$; if
$W^{(\mathrm{loc})}_{jv}$ is large, location $v$ belongs to group $j$. The core
$w^{(\mathrm{core})}_{ij}(t)$ is the activity curve of that group pair.

\begin{keypoint}
Equation~\eqref{eq:trend-elementwise} is the heart of interpretability: the value for $(u,v)$ is a sum
of latent dynamics, each weighted by how strongly the keyword and location relate to it. Non-negative
weights keep the sum from flipping direction.
\end{keypoint}

\subsection{The core is generated by reaction--diffusion}

The paper does not let $w^{(\mathrm{core})}_{ij}(t)$ vary arbitrarily. It assumes the core is generated
by the system in Equation~\eqref{eq:rd-latent} from Section~\ref{sec:reaction-diffusion}, paper Eq.~3:
\begin{equation}
  \frac{dw_{ij}}{dt}
  =a_{ij}w_{ij}
   +\sum_{j'=1}^{d_l}d_{ijj'}(w_{ij'}-w_{ij}).
  \label{eq:core-ode}
\end{equation}
This is a system of \textbf{ordinary differential equations (ODEs)} in time. Data are observed at
discrete steps (weeks or days), but the model defines continuous-time dynamics and integrates them to
produce a sequence of $L_c$ steps. Positions are discrete (location groups); time in the equation is
continuous.

Let $\theta=\{\mathbf{A},\mathcal{D}\}$ be the system parameters, and let
$w_0\in\mathbb{R}_{+}^{d_k\times d_l}$ be the initial state. The paper writes
$f_\theta(\theta,w_0,L)$ for the function that generates an $L$-long core sequence from $w_0$. The full
trend parameters are therefore
\[
  \theta_d=\{\mathbf{W}^{(\mathrm{key})},\mathbf{W}^{(\mathrm{loc})},\theta\}.
\]
In words, $\theta_d$ contains both ``how the dynamics evolve'' and ``which keywords/locations they
attach to''.

\subsection{Seasonal: a similar factorization}

The seasonal part is compressed too, but without an ODE. The paper factorizes (Eq.~4):
\begin{equation}
  \widehat{\mathcal{X}}_s
  =\mathbf{S}^{(\mathrm{time})}
   \times_{\mathrm{key}}\mathbf{S}^{(\mathrm{key})}
   \times_{\mathrm{loc}}\mathbf{S}^{(\mathrm{loc})},
  \label{eq:seasonal-factorization}
\end{equation}
with $d_s$ seasonal patterns and
\[
  \mathbf{S}^{(\mathrm{time})}\in\mathbb{R}^{d_s\times L_c},
  \quad
  \mathbf{S}^{(\mathrm{key})}\in\mathbb{R}^{d_s\times k},
  \quad
  \mathbf{S}^{(\mathrm{loc})}\in\mathbb{R}^{d_s\times l}.
\]
Elementwise,
\begin{equation}
  \widehat{x}_{s,uv}(t)
  =\sum_{r=1}^{d_s}
    S^{(\mathrm{time})}_{r,t}S^{(\mathrm{key})}_{r,u}S^{(\mathrm{loc})}_{r,v}.
  \label{eq:seasonal-elementwise}
\end{equation}
Each $r$ is a seasonal ``template'' over time; the other factors say how strong it is for each
keyword and location. This is a PARAFAC factorization: each pattern is a rank-one sum of three
vectors. The paper does not state that seasonal factors must be non-negative like the trend
$\mathbf{W}$, so this guide does not add that constraint.

\subsection{Outliers: sparse, not all noise}

The outlier tensor $\widehat{\mathcal{X}}_o$ is assumed \emph{sparse}: most entries are zero, retaining
only unusual values not explained by trend and seasonality. During optimization, the paper takes the
residual $\mathcal{X}^{c}-\widehat{\mathcal{X}}_d-\widehat{\mathcal{X}}_s$ and keeps entries only
when they reduce the total MDL cost. ``Outlier'' is therefore a structured component, not a label for
``bad data''.

\subsection{Parameter sets and multiple regimes}

A single model is packaged as
\[
  \Theta=\{\theta_d,\theta_s,\widehat{\mathcal{X}}_o\},
  \qquad
  \theta_s=\{\mathbf{S}^{(\mathrm{time})},\mathbf{S}^{(\mathrm{key})},
             \mathbf{S}^{(\mathrm{loc})}\}.
\]
When patterns change, one $\Theta$ is insufficient. The paper allows $N$ models:
\[
  \mathcal{F}=\{\Theta_1,\Theta_2,\ldots,\Theta_N\}.
\]
Each $\Theta_n$ describes a relatively stable period; $N$ and switch times are selected automatically
by MDL in Section~\ref{sec:mdl-streaming}, not split by eye. Thus, in a plot such as
Figure~\ref{fig:preview}(a), the dotted lines are not data events; they are model boundaries.

\begin{center}
\small
\begin{tabular}{@{}p{0.24\linewidth}p{0.68\linewidth}@{}}
  \toprule
  Parameter & What it contains \\
  \midrule
  $\theta$ & $\mathbf{A}$ (growth) and $\mathcal{D}$ (diffusion) of the latent system. \\
  $\theta_d$ & Trend core, keyword/location factors, and $\theta$. \\
  $\theta_s$ & Seasonal patterns and keyword/location factors. \\
  $\Theta$ & One complete model: trend, seasonality, and outliers. \\
  $\mathcal{F}$ & A collection of $\Theta$ models for different stream periods. \\
  \bottomrule
\end{tabular}
\end{center}

\subsection{Reading one dynamics from core to raw data}

Suppose the model learns $d_k=3$ keyword groups and $d_l=2$ location groups. To understand
``Dynamics \#1 of Group 1'' in Figure~\ref{fig:preview}(b), follow four steps:
\begin{enumerate}
  \item Look at the core curve: is the dynamics rising, falling, or changing direction?
  \item Look at the corresponding keyword-factor column: which keywords have high weights? In the
        2019--2021 example, the paper associates ``iPhone'' strongly with Dynamics \#1.
  \item Look at the location-factor row: which groups have high weights? The paper describes Group~1
        as the United States, Europe, and China, and Group~2 as the rest of the world.
  \item Inspect $d_{ijj'}$ and the sign of $(w_{ij'}-w_{ij})$ to read the direction. The paper says
        Dynamics \#1 spreads from Group~1 to Group~2, i.e.\ interest in ``iPhone'' diffused from the
        leading regions to the rest of the world in that period.
\end{enumerate}
This is \emph{model-based interpretation}, not a human-provided label. The name ``Group 1'' only gains
meaning after the factors are learned; it need not match a continent or language.

\begin{plainwords}
To read the result, picture three layers: the small core says ``what is happening'', the factor
rows/columns say ``where and for which keyword'', and the diffusion arrows say ``from which group to
which group''. You do not need to stare at the raw tensor.
\end{plainwords}

\begin{commonmistake}
A common error is to treat $d_l\ll k$ in the paper's prose as a real constraint. The shape of
$\mathbf{W}^{(\mathrm{loc})}$ is $d_l\times l$, and the intent is to reduce the number of location
groups, so this guide uses $d_k\ll k$, $d_l\ll l$. This is a likely typo in the paper, surfaced rather
than copied silently. Another error is to conflate $\Theta$ and $\mathcal{F}$: $\mathcal{F}$ is a
\emph{set} of models for different periods.
\end{commonmistake}

\begin{selfcheck}
\begin{enumerate}
  \item Write the shapes of the four objects in Equation~\eqref{eq:trend-factorization}.
  \item Distinguish $\theta$, $\theta_d$, $\Theta$, and $\mathcal{F}$ in one sentence each.
  \item Why does seasonality not need a reaction--diffusion ODE? Which assumption differs from trend?
  \item If $N=1$ for the whole stream, what is the model assuming?
\end{enumerate}
\end{selfcheck}
